\section{Shared Resources}

\subsection{Race Conditions}

\begin{definition}
A \textbf{Race condition} occurs in concurent programming when the correctness of the system depends on
the specific interleaving or ordering of execution of operations.
\end{definition}

\begin{definition}
A \textbf{Low-level race condition} is a race condition that occurs because of insufficient
synchronization of memory access (Data Races). While the \textbf{High-level race condition} is a race condition that
occurs because of a wrong decomposition/algorithm. The program might ensure mutual exclusion on the
shared resources but bad interleavings still lead to wrong results.
\end{definition}

Generally we want one of three things when it comes to shared ressources:
\begin{itemize}
    \item \textbf{Thread local (Isolated Mutability):} We dont use the ressource of other threads.
    \item \textbf{Read only (Immutable):} We only read and do not change the ressource
    \item \textbf{Mutual exclusion (Synchronization):} We use the ressource of other threads but we ensure that only one thread accesses it at a time.
\end{itemize}

\subsection{Guidelines for Concurrent programming}

\begin{important}
\textbf{Thread-Locality:} Minimize the sharing of mutable state between threads. Only share ressources if needed.
\end{important}

\begin{important}
\textbf{Immutability:} Use immutable data structures whenever possible. They are inherently thread-safe.
\end{important}

\begin{important}
\textbf{Consitent locking:} If you need to share mutable state, use locks to ensure mutual exclusion.
\end{important}

\begin{definition}
\textbf{Lock granularity } refers to the amount of code that is proteted by a lock in ways:
\begin{itemize}
        \item Size of synchronized blocks
        \item Number of objects a lock is used for
\end{itemize}
\end{definition}

\begin{theorem}
\textbf{Coarse-grained locking} is easier to implement but leads to less concurrency and therfore performance loss. \textbf{Fine-grained locking} is harder to implement without creating bugs but leads to more concurrency and therfore performance gain.
\end{theorem}

\begin{important}
Start with coarse-grained (simpler) and move to fine-
grained (performance) only if contention on the coarser locks becomes an issue.
\end{important}

\begin{important}
Do not do expensive computations or I/O in critical
sections, but also don’t introduce race conditions
\end{important}

\begin{definition}
\textbf{Atomicity: } An operation is atomic if it is performed as a single, indivisible unit. In other words, it either completes entirely or not at all, without any intermediate states being visible to other threads.
\end{definition}


\begin{important}
Think in terms of what operations need to be atomic
\end{important}